\section*{Hoofdstuk \ref{ch:b1:titration-methods} --- Titratiemethoden en titratiekrommen}

\begin{solution}{exo:b1:titration-methods:1}
\ce{H3O+ + OH- -> 2H2O}, $K = 1/K_e = 10^{14.00}$. $C = 12.4 \times
0.100/20.0 = \qty{0.0620}{mol/L}$.
\end{solution}


\begin{solution}{exo:b1:titration-methods:2}
Bij de \omterm{def:b1:titration-methods:titration}{equivalentie} is de oplossing ammoniumchloride van
\qty{0.050}{mol/L}: $\mathrm{pH} = \frac12(9.25 - \log 0.050) = 5.28$.
Methylrood (3.8--5.8) omvat die waarde; methyloranje, waarvan het traject
(2.5--4.5) pas na de \omterm{def:b1:titration-methods:titration}{equivalentie} bereikt wordt, wanneer de pH daalt, zou
te laat omslaan.
\end{solution}

\begin{solution}{exo:b1:titration-methods:3}
\ce{MnO4- + 5Fe^2+ + 8H+ -> Mn^2+ + 5Fe^3+ + 4H2O};
$C = 5 \times 0.0200 \times 12.6/10.0 = \qty{0.126}{mol/L}$.
\end{solution}

\begin{solution}{exo:b1:titration-methods:4}
$C = 0.0100 \times 14.2/50.0 = \qty{2.84e-3}{mol/L}$ calcium plus
magnesium. Beide reageren vóór het \omterm{def:b1:titration-methods:titration}{eindpunt}, dat dus alleen hun som
geeft: een \omterm{def:b1:titration-methods:kinds}{gelijktijdige titratie}.
\end{solution}

\begin{solution}{exo:b1:titration-methods:5}
0~mL: \omterm{def:b1:predominant-reaction:strong-acid}{zwak zuur}, $\mathrm{pH} = \frac12(4.76 + 1.00) = 2.88$. 5.0~mL:
halve \omterm{def:b1:titration-methods:titration}{equivalentie}, 4.76. 10.0~mL: 8.73
(\cref{prop:b1:titration-methods:ph-at-equivalence}). 15.0~mL: overmaat
hydroxide \qty{0.50}{mmol} in \qty{25.0}{mL}, $[\ce{OH-}] = 0.020$,
$\mathrm{pH} = 12.30$. Alles stemt tot op 0.01 overeen met de berekende
kromme.
\end{solution}

\begin{solution}{exo:b1:titration-methods:6}
De $\mathrm{p}K_a$-waarden 2.15 en 7.21 verschillen meer dan 4, net als
7.21 en 12.34: twee \omterm{def:b1:titration-methods:titration}{equivalenties}, bij 10.0 en \qty{20.0}{mL}. Eerste:
oplossing van \ce{H2PO4-}, $\mathrm{pH} \approx \frac12(2.15 + 7.21) = 4.68$;
tweede: \ce{HPO4^2-}, $\frac12(7.21 + 12.34) = 9.78$ (de exacte kromme
geeft 4.71 en 9.66, omdat de formules verdunning en water verwaarlozen).
De derde zuurfunctie ($\mathrm{p}K_a$ 12.34) ligt te dicht bij die van
water: het toegevoegde hydroxide reageert maar gedeeltelijk en de pH
stijgt zonder sprong.
\end{solution}

\begin{solution}{exo:b1:titration-methods:7}
Zoutzuur: $6.2 \times 0.100/10.0 = \qty{0.062}{mol/L}$; ethaanzuur:
$(14.6 - 6.2) \times 0.100/10.0 = \qty{0.084}{mol/L}$. Halverwege is de
helft van het ethaanzuur getitreerd: $\mathrm{pH} = \mathrm{p}K_a = 4.76$.
\end{solution}

\begin{solution}{exo:b1:titration-methods:8}
$V_{\mathrm{eq}} = \qty{10.0}{mL}$. 2.0~mL: $E = 0.77 + 0.059\log(2/8) =
0.73$~V; 9.0~mL: $0.77 + 0.059\log 9 = 0.83$~V; 11.0~mL: overmaat
permanganaat \qty{0.020}{mmol}, \ce{Mn^2+} \qty{0.200}{mmol}, $E = 1.51 +
\frac{0.059}{5}\log 0.10 = 1.50$~V; 20.0~mL: 1.51~V.
\end{solution}

\begin{solution}{exo:b1:titration-methods:9}
0~mL: $\kappa = (349.8 + 76.2) \times 0.0100 = \qty{4.26}{mS/cm}$.
10.0~mL: \ce{Na+} en \ce{Cl-} op $1.00/110 = \qty{9.09e-3}{mol/L}$,
$\kappa = (50.3 + 76.2) \times \num{9.09e-3} = 1.15$. 20.0~mL: \ce{Na+}
$0.0167$, \ce{Cl-} en \ce{OH-} $0.00833$~mol/L, $\kappa = 0.84 + 0.635 + 1.65
= 3.13$~mS/cm. Hellingen (verdunning verwaarloosd, \qty{1}{mL} voegt
\qty{1.0e-3}{mol/L} toe): $(50.3 - 349.8) \times 10^{-3} = -0.30$ en
$(50.3 + 198.3) \times 10^{-3} = +0.25$~mS/cm per mL.
\end{solution}

\begin{solution}{exo:b1:titration-methods:10}
Bij $V_{\mathrm{eq}} - v$ is het overmaat zuur $C_Bv$, bij
$V_{\mathrm{eq}} + v$ is de overmaat base $C_Bv$; in hetzelfde volume is
$[\ce{H3O+}]$ ervoor gelijk aan $[\ce{OH-}]$ erna, dus
$\mathrm{pH}(V_{\mathrm{eq}} - v) +
\mathrm{pH}(V_{\mathrm{eq}} + v) = 14$: symmetrie rond (V$_{\mathrm{eq}}$,
7). Sprong: \qty{1.0e-5}{mol} overmaat in \qty{20}{mL}, pH 3.30 tot
10.70, 7.4 eenheden. Gedeeld door 100: $\num{1.0e-7}$~mol in
\qty{20}{mL}, pH 5.3 tot 8.7, 3.4 eenheden: een veel minder scherp
\omterm{def:b1:titration-methods:titration}{eindpunt}.
\end{solution}

\begin{solution}{exo:b1:titration-methods:11}
Ervoor: gevormd \ce{Fe^3+} $= 5C_{\mathrm{Mn}}V$, overgebleven \ce{Fe^2+} $=
5C_{\mathrm{Mn}}(V_{\mathrm{eq}} - V)$, dus $E = 0.77 + 0.059\log\frac{V}{V_{\mathrm{eq}}
- V}$: bij $V_{\mathrm{eq}}/2$ is $E = 0.77$. Erna: overmaat \ce{MnO4-} $\propto
V - V_{\mathrm{eq}}$, \ce{Mn^2+} $\propto V_{\mathrm{eq}}$: $E = 1.51 +
\frac{0.059}{5}\log\frac{V - V_{\mathrm{eq}}}{V_{\mathrm{eq}}}$ (pH 0); bij
$2V_{\mathrm{eq}}$ 1.51~V. Bij de \omterm{def:b1:titration-methods:titration}{equivalentie} $(0.77 + 5 \times 1.51)/6 =
1.39$~V. Bij pH 1 ligt het permanganaatkoppel $0.059 \times 8/5 =
0.094$~V lager: $E_{\mathrm{eq}} = (0.77 + 5 \times 1.416)/6 = 1.31$~V.
\end{solution}

\begin{solution}{exo:b1:titration-methods:12}
\ce{NH4+ + OH- -> NH3 + H2O} heeft $K = 10^{14.00 - 9.25} = 10^{4.75}$:
de reactie is rond de \omterm{def:b1:titration-methods:titration}{equivalentie} niet volledig en de sprong is klein,
rond pH 11 (pH bij de \omterm{def:b1:titration-methods:titration}{equivalentie} $14 - \frac12(4.75 + 1.30) = 11.0$).
\omterm{def:b1:titration-methods:kinds}{Terugtitratie}: voeg $n_1$ natronloog toe (overmaat), kook om de gevormde
ammoniak te verdrijven, en titreer daarna het overgebleven hydroxide met
zoutzuur ($n_2$, scherpe sprong sterk--sterk): $n(\ce{NH4+}) = n_1 - n_2$.
\end{solution}

\begin{solution}{pb:b1:titration-methods:1}
\textbf{1.} \ce{C6H6O6 + 2H+ + 2e- -> C6H8O6}.
\textbf{2.} \ce{C6H8O6 + I2 -> C6H6O6 + 2I- + 2H+}.
\textbf{3.} 0 en $-$I.
\textbf{4.} \ce{I2 + 2S2O3^2- -> 2I- + S4O6^2-}, $\log K = 2(0.53 - 0.02)/0.059
= 17$.
\textbf{5.} De blauwe kleur vereist jood; ze verdwijnt met het laatste
jood, bij de \omterm{def:b1:titration-methods:titration}{equivalentie}.
\textbf{6.} Verliezen aan de lucht zouden het resultaat verlagen; meteen
en in overmaat toegevoegd oxideert het jood het ascorbinezuur snel en
volledig.
\textbf{7.} Een snelle, \omterm{def:b1:extent-q-and-k:quantitative}{kwantitatieve reactie} gedurende de hele
\omterm{def:b1:titration-methods:titration}{titratie}, zonder oxidatie door de lucht terwijl de \omterm{def:b1:titration-methods:titration}{titrant} langzaam
toegevoegd wordt, en een \omterm{def:b1:titration-methods:titration}{eindpunt} bij de eerste overmaat jood.
\textbf{8.} Overmaat jood $n_R$; reactie met ascorbinezuur; \omterm{def:b1:titration-methods:titration}{titratie} van
het overgebleven jood met thiosulfaat.
\textbf{9.} Anders zou er ascorbinezuur overblijven, en zou de overmaat
jood, nul, niet zeggen hoeveel. Hier heeft
\qty{2.775e-4}{mol} gereageerd van de \qty{5.000e-4}{mol}: inderdaad een
overmaat.
\textbf{10.} Thiosulfaatoplossingen veranderen langzaam in de tijd; de
concentratie ervan komt rechtstreeks in het resultaat terecht.
\textbf{11.} 10.
\textbf{12.} Een volpipet van \qty{20.00}{mL}.
\textbf{13.} $20.00 \times 10^{-3} \times 0.0250 = \qty{5.000e-4}{mol}$.
\textbf{14.} $8.90 \times 10^{-3} \times 0.0500 = \qty{4.450e-4}{mol}$.
\textbf{15.} De helft: \qty{2.225e-4}{mol}.
\textbf{16.} $5.000 - 2.225 = \qty{2.775e-4}{mol}$.
\textbf{17.} Eén op één: \qty{2.775e-4}{mol}.
\textbf{18.} $\times 10$: \qty{2.775e-3}{mol}, $\times 176.0$:
\qty{0.488}{g}.
\textbf{19.} Ongeveer \qty{2}{\percent} onder het etiket.
\textbf{20.} $n = \bigl(C_IV_I - \frac12C_TV_T\bigr)\,V_{\mathrm{kolf}}/V_{\mathrm{deel}}$.
\textbf{21.} Ingebracht jood: pipet $0.03/20.00 = \qty{0.15}{\percent}$
en concentratie \qty{0.2}{\percent}, ongeveer \qty{0.25}{\percent}, dat
is \qty{1.3e-6}{mol}. Overmaat: buret $0.05/8.90 = \qty{0.56}{\percent}$
en concentratie \qty{0.2}{\percent}, ongeveer \qty{0.60}{\percent}, dat
is \qty{1.3e-6}{mol}.
\textbf{22.} De absolute onzekerheden tellen samen, terwijl het verschil
kleiner is dan elk van beide hoeveelheden: $\sqrt{1.3^2 + 1.3^2} \times 10^{-6} =
\qty{1.8e-6}{mol}$, \qty{0.66}{\percent} van \qty{2.775e-4}{mol}, meer
dan elk van beide termen.
\textbf{23.} Met de maatkolf (\qty{0.1}{\percent}) en de pipet van het
genomen deel (\qty{0.2}{\percent}):
$\sqrt{0.66^2 + 0.1^2 + 0.2^2} = \qty{0.70}{\percent}$.
\textbf{24.} $0.488 \pm 0.003$~g: \textbf{\qty{0.49}{g} ascorbinezuur
per tablet}.
\end{solution}
